% ===========================================================================
%  Jàng Afrig — Lab 3 : CSS en pratique
%  Styliser les 5 pages avec palette panafricaine + Google Fonts
%  Licence 1 DUT1-Informatique — ESP/UCAD — 2025–2026
%  Dr. El Hadji Bassirou TOURÉ
%  Compilé avec : xelatex (2–3 passes)
% ===========================================================================
\documentclass[11pt,a4paper]{article}

\usepackage{fontspec}
\setmainfont{Latin Modern Roman}
\setsansfont{Latin Modern Sans}
\setmonofont{Latin Modern Mono}
\usepackage{polyglossia}
\setdefaultlanguage{french}

\usepackage[margin=2.5cm]{geometry}
\usepackage{amsmath,amssymb}
\usepackage{graphicx}
\graphicspath{{figures/}}
\usepackage{xcolor}
\usepackage{booktabs}
\usepackage{hyperref}
\usepackage{tcolorbox}
\tcbuselibrary{theorems,breakable,skins}
\usepackage{enumitem}
\usepackage{fancyhdr}
\usepackage{array}
\usepackage{pifont}
\usepackage{longtable}
\usepackage{multirow}

% --- TikZ ---
\usepackage{tikz}
\usetikzlibrary{shapes.geometric,arrows.meta,positioning,calc,backgrounds,fit,shadows.blur}

\newcommand{\cmark}{\ding{51}}
\newcommand{\xmark}{\ding{55}}

\hypersetup{
    colorlinks=true,
    linkcolor=blue!60!black,
    urlcolor=blue!60!black,
    citecolor=blue!60!black
}

\newtcolorbox{objectifbox}[1][]{colback=blue!5,colframe=blue!50,fonttitle=\bfseries,title=#1,breakable}
\newtcolorbox{infobox}[1][]{colback=green!5,colframe=green!50!black,fonttitle=\bfseries,title=#1,breakable}
\newtcolorbox{attentionbox}[1][]{colback=orange!5,colframe=orange!70,fonttitle=\bfseries,title=#1,breakable}
\newtcolorbox{retenirbox}[1][]{colback=yellow!10,colframe=yellow!50!black,fonttitle=\bfseries,title=#1,breakable}
\newtcolorbox{prereqbox}[1][]{colback=cyan!5,colframe=cyan!50!black,fonttitle=\bfseries,title=#1,breakable}
\newtcolorbox{projetbox}[1][]{colback=teal!5,colframe=teal!50!black,fonttitle=\bfseries,title=#1,breakable}

% --- Couleurs palette panafricaine ---
\definecolor{panvert}{HTML}{0E8A4F}
\definecolor{panjaune}{HTML}{F5C518}
\definecolor{panrouge}{HTML}{C8332E}
\definecolor{pannoir}{HTML}{1A1A1A}
\definecolor{panblanc}{HTML}{F8F8F8}

% --- Code source ---
\usepackage{listings}
\lstdefinelanguage{HTML5}{
  morekeywords={<!DOCTYPE,html,head,meta,title,body,h1,h2,h3,h4,p,ul,ol,li,table,tr,th,td,a,img,strong,em,br,div,span,header,nav,main,section,article,footer,aside,form,input,select,option,textarea,label,button,fieldset,legend,style,link,class,id,href,src,alt},
  sensitive=true,
  morecomment=[s]{<!--}{-->},
  morestring=[b]",
}
\lstdefinelanguage{CSS}{
  morekeywords={color,background-color,background,font-family,font-size,font-weight,text-align,padding,margin,border,border-radius,width,height,max-width,min-width,display,flex,flex-direction,justify-content,align-items,flex-wrap,gap,box-sizing,transition,transform,position,top,bottom,left,right,opacity,cursor,line-height,letter-spacing,box-shadow,list-style,text-decoration,text-transform,border-collapse,vertical-align,object-fit,overflow},
  sensitive=true,
  morecomment=[s]{/*}{*/},
  morestring=[b]",
}
\lstset{
  basicstyle=\ttfamily\small,
  backgroundcolor=\color{gray!8},
  frame=single,
  framerule=0.5pt,
  rulecolor=\color{gray!40},
  breaklines=true,
  breakatwhitespace=true,
  showstringspaces=false,
  tabsize=2,
  xleftmargin=4pt,
  xrightmargin=4pt,
  aboveskip=8pt,
  belowskip=8pt,
  literate={é}{{\'e}}1 {è}{{\`e}}1 {ê}{{\^e}}1 {à}{{\`a}}1 {ù}{{\`u}}1
           {ô}{{\^o}}1 {î}{{\^i}}1 {ç}{{\c{c}}}1 {É}{{\'E}}1
           {→}{{$\rightarrow$}}1 {≈}{{$\approx$}}1
}

\pagestyle{fancy}
\fancyhf{}
\fancyhead[L]{\small Jàng Afrig --- Lab 3 : CSS en pratique}
\fancyhead[R]{\small UCAD --- L1 DUT1-Info}
\fancyfoot[C]{\thepage}
\fancyfoot[L]{\small Dr. El Hadji Bassirou TOURÉ}
\fancyfoot[R]{\small 2025--2026}
\renewcommand{\headrulewidth}{0.4pt}
\renewcommand{\footrulewidth}{0.4pt}

\title{%
\vspace{-1cm}
\textbf{\LARGE Fondamentaux du Web}\\
\textbf{\LARGE HTML, CSS, JavaScript}\\[0.8cm]
\textbf{\huge Lab 3 --- CSS en pratique}\\[0.3cm]
\textbf{\Large Styliser Jàng Afrig avec une palette panafricaine}\\[0.5cm]
\large Licence 1 DUT1-Informatique --- 2025--2026}
\author{Dr. El Hadji Bassirou TOURÉ\\
Département Génie Informatique\\
Université Cheikh Anta Diop de Dakar}
\date{}

\begin{document}

\maketitle
\thispagestyle{empty}

\vspace{0.5cm}

% ===========================================================================
% BOÎTE PRÉREQUIS
% ===========================================================================
\begin{prereqbox}[Informations pratiques]
\begin{center}
\begin{tabular}{ll}
\toprule
\textbf{Type} & Lab guidé \\
\textbf{Séance} & S5 --- premier Lab CSS \\
\textbf{Prérequis} & Lab 2 terminé (5 pages HTML poussées avec tag \texttt{v0.2}) + CM2 \\
\textbf{Validation} & 5 pages stylisées + tag \texttt{v0.3} sur GitHub \\
\textbf{Rapport} & 1-2 pages : ce que j'ai fait + difficultés + lien repo \\
\bottomrule
\end{tabular}
\end{center}
\end{prereqbox}

\begin{objectifbox}[Objectifs du Lab]
À la fin de ce Lab, vous aurez :
\begin{enumerate}[nosep]
    \item Créé un fichier \texttt{style.css} et l'avoir lié aux 5 pages HTML
    \item Mis en place un \textbf{reset CSS} et \texttt{box-sizing: border-box}
    \item Importé \textbf{2 polices Google Fonts} (titres + corps)
    \item Défini une \textbf{palette panafricaine} (vert, jaune, rouge + neutres)
    \item Stylisé le \textbf{header} et la \textbf{navigation Flexbox}
    \item Créé des \textbf{cartes} pour les modes de jeu
    \item Stylisé des \textbf{boutons} avec effet \texttt{:hover}
    \item Stylisé le \textbf{tableau} du palmarès (bordures + lignes alternées)
    \item Posé le tag \textbf{\texttt{v0.3}} sur GitHub
\end{enumerate}
\end{objectifbox}

\tableofcontents

\newpage

% ===========================================================================
\section{Avant de coder}
% ===========================================================================

\subsection{Là où on en est}

Au Lab 2, vous avez construit les 5 pages HTML de Jàng Afrig avec une
structure sémantique propre. Mais visuellement, c'est encore austère :
texte noir sur fond blanc, liens bleus soulignés, tableau sans bordures.

\subsection{Ce qu'on va construire aujourd'hui}

Un \textbf{fichier \texttt{style.css}} unique, lié aux 5 pages, qui les
transforme en application visuellement attractive.

\begin{figure}[!ht]
\centering
\begin{tikzpicture}[
    every node/.style={anchor=north west, inner sep=8pt, font=\scriptsize}
]
    \draw[draw=gray!50, thick, rounded corners=4pt, fill=panblanc] (0,0) rectangle (12,9);

    % Header panvert
    \filldraw[fill=panvert, draw=panvert] (0.05, 8.4) rectangle (11.95, 8.95);
    \node[font=\small\bfseries, color=white, anchor=center] at (6, 8.7) {Jàng Afrig};

    % Nav panjaune
    \filldraw[fill=panjaune!50, draw=panjaune!70!black] (0.05, 7.8) rectangle (11.95, 8.4);
    \node[font=\small, anchor=center] at (3, 8.1) {Accueil};
    \node[font=\small, anchor=center] at (5.5, 8.1) {Jouer};
    \node[font=\small, anchor=center] at (8, 8.1) {Mes Scores};
    \node[font=\small, anchor=center] at (10.5, 8.1) {À propos};

    % Section Bienvenue
    \node[font=\bfseries\small, color=panvert] at (0.5, 7.5) {Bienvenue};
    \node[anchor=north west] at (0.5, 7.05) {Bienvenue sur Jàng Afrig, un quiz géographique sur les 54 pays africains.};

    % Section Modes de jeu (cartes)
    \node[font=\bfseries\small, color=panvert] at (0.5, 6.2) {Les modes de jeu};

    \filldraw[fill=panjaune!20, draw=panjaune!70!black, thick, rounded corners=4pt] (0.5, 4.3) rectangle (3.5, 5.7);
    \node[font=\small\bfseries, anchor=north] at (2, 5.55) {Capitales};
    \node[anchor=north, font=\tiny, text width=2.6cm, align=center] at (2, 5.1) {Quelle est la capitale du Mali ?};

    \filldraw[fill=panrouge!20, draw=panrouge!70!black, thick, rounded corners=4pt] (3.7, 4.3) rectangle (6.7, 5.7);
    \node[font=\small\bfseries, anchor=north] at (5.2, 5.55) {Drapeaux};
    \node[anchor=north, font=\tiny, text width=2.6cm, align=center] at (5.2, 5.1) {À quel pays appartient ce drapeau ?};

    \filldraw[fill=panvert!20, draw=panvert!70!black, thick, rounded corners=4pt] (6.9, 4.3) rectangle (9.9, 5.7);
    \node[font=\small\bfseries, anchor=north] at (8.4, 5.55) {Langues};
    \node[anchor=north, font=\tiny, text width=2.6cm, align=center] at (8.4, 5.1) {Dans quel pays parle-t-on swahili ?};

    % Bouton CTA
    \filldraw[fill=panvert, draw=panvert!70!black, thick, rounded corners=4pt] (5, 3.3) rectangle (7, 3.85);
    \node[font=\small\bfseries, color=white, anchor=center] at (6, 3.575) {Commencer};

    % Tableau palmarès
    \node[font=\bfseries\small, color=panvert] at (0.5, 3) {Palmarès};
    \filldraw[fill=panvert!10, draw=panvert!70!black, thick] (0.5, 2.5) rectangle (11.5, 2.85);
    \node[font=\tiny, anchor=west] at (0.8, 2.675) {Rang};
    \node[font=\tiny, anchor=west] at (3, 2.675) {Pseudo};
    \node[font=\tiny, anchor=west] at (5.5, 2.675) {Mode};
    \node[font=\tiny, anchor=west] at (8, 2.675) {Score};
    \filldraw[fill=panblanc, draw=panvert!40, thick] (0.5, 2.15) rectangle (11.5, 2.5);
    \node[font=\tiny, anchor=west] at (0.8, 2.325) {1};
    \node[font=\tiny, anchor=west] at (3, 2.325) {Fatou};
    \node[font=\tiny, anchor=west] at (5.5, 2.325) {Capitales};
    \node[font=\tiny, anchor=west] at (8, 2.325) {10/10};
    \filldraw[fill=panvert!5, draw=panvert!40, thick] (0.5, 1.8) rectangle (11.5, 2.15);
    \node[font=\tiny, anchor=west] at (0.8, 1.975) {2};
    \node[font=\tiny, anchor=west] at (3, 1.975) {Moussa};
    \node[font=\tiny, anchor=west] at (5.5, 1.975) {Drapeaux};
    \node[font=\tiny, anchor=west] at (8, 1.975) {9/10};

    % Footer
    \filldraw[fill=pannoir, draw=pannoir] (0.05, 0.05) rectangle (11.95, 0.7);
    \node[font=\tiny, color=white, anchor=center] at (6, 0.375) {© 2025--2026 Jàng Afrig --- L1 DUT1-Info, ESP/UCAD};
\end{tikzpicture}
\caption{Aperçu de l'objectif visuel à atteindre dans ce Lab.}
\label{fig:apercu}
\end{figure}

\textbf{Header vert, nav jaune, cartes colorées (vert/jaune/rouge), bouton
vert, palmarès avec lignes alternées, footer noir.} On y va.

\newpage

% ===========================================================================
\section{Étape 1 : Créer style.css et le lier aux 5 pages}
% ===========================================================================

\subsection{Créer le fichier}

Dans VS Code, à la racine de votre repo \texttt{jang-afrig} :
\begin{enumerate}[nosep]
    \item Clic sur l'icône «\,Nouveau fichier\,» dans l'Explorateur
    \item Nommez-le : \texttt{style.css}
    \item Laissez-le vide pour l'instant (on va le remplir étape par étape)
\end{enumerate}

\subsection{Lier style.css aux 5 pages}

Dans \textbf{chaque} fichier HTML (\texttt{index.html}, \texttt{quiz.html},
\texttt{resultat.html}, \texttt{scores.html}, \texttt{apropos.html}),
ajoutez cette ligne dans le \texttt{<head>} :

\begin{lstlisting}[language=HTML5]
<head>
    <meta charset="UTF-8">
    <title>Jàng Afrig</title>
    <link rel="stylesheet" href="style.css">
</head>
\end{lstlisting}

\textbf{Enregistrez les 5 fichiers}. Pour l'instant, comme \texttt{style.css}
est vide, rien ne change visuellement. C'est attendu.

\begin{retenirbox}[Le piège du copier-coller]
Si vous oubliez d'ajouter le \texttt{<link>} sur \emph{une seule} page,
celle-ci restera moche pendant que les autres seront stylées. Vérifiez
bien les 5.
\end{retenirbox}

% ===========================================================================
\section{Étape 2 : Reset CSS et box-sizing}
% ===========================================================================

Chaque navigateur applique des styles par défaut différents. Pour démarrer
sur une base propre, on ajoute un \emph{reset minimal}.

\textbf{Tapez ceci au début de \texttt{style.css}} :

\begin{lstlisting}[language=CSS]
/* ===== Reset minimal ===== */
* {
    box-sizing: border-box;
    margin: 0;
    padding: 0;
}

body {
    font-family: sans-serif;
    line-height: 1.6;
    color: #1A1A1A;
    background-color: #F8F8F8;
}
\end{lstlisting}

\textbf{Enregistrez}. La page change d'apparence : tous les espacements
sont remis à zéro, les marges des paragraphes disparaissent. C'est moche
provisoirement, mais on repart d'une base contrôlée. \cmark

\begin{infobox}[Pourquoi ce reset ?]
\begin{itemize}[nosep]
    \item \texttt{box-sizing: border-box} : les calculs de largeur incluent padding et border (vu au CM2)
    \item \texttt{margin: 0; padding: 0} : on annule les marges par défaut du navigateur
    \item \texttt{line-height: 1.6} : un peu d'air entre les lignes pour la lisibilité
\end{itemize}
\end{infobox}

\newpage

% ===========================================================================
\section{Étape 3 : La palette panafricaine}
% ===========================================================================

\subsection{Le choix des couleurs}

On utilise les couleurs panafricaines (présentes sur de nombreux drapeaux
du continent) :

\begin{center}
\begin{tabular}{@{}lll@{}}
\toprule
\textbf{Couleur} & \textbf{Code hex} & \textbf{Usage} \\
\midrule
\textcolor{panvert}{$\blacksquare$} Vert  & \texttt{\#0E8A4F} & Header, accents, bouton principal \\
\textcolor{panjaune}{$\blacksquare$} Jaune & \texttt{\#F5C518} & Nav, mise en valeur \\
\textcolor{panrouge}{$\blacksquare$} Rouge & \texttt{\#C8332E} & Alertes, mode «\,Drapeaux\,» \\
\textcolor{pannoir}{$\blacksquare$} Noir   & \texttt{\#1A1A1A} & Texte principal, footer \\
\textcolor{black}{$\square$} Blanc cassé   & \texttt{\#F8F8F8} & Fond de page \\
\bottomrule
\end{tabular}
\end{center}

\subsection{Variables CSS : un seul endroit pour les couleurs}

Plutôt que de répéter \texttt{\#0E8A4F} 20 fois dans le fichier, on définit
des \textbf{variables CSS} (officiellement appelées \emph{custom properties}).

\textbf{Ajoutez en haut de \texttt{style.css}, juste après le reset :}

\begin{lstlisting}[language=CSS]
/* ===== Palette panafricaine ===== */
:root {
    --vert: #0E8A4F;
    --jaune: #F5C518;
    --rouge: #C8332E;
    --noir: #1A1A1A;
    --blanc-casse: #F8F8F8;
}
\end{lstlisting}

\textbf{Comment les utiliser} dans n'importe quelle règle :

\begin{lstlisting}[language=CSS]
header {
    background-color: var(--vert);
}
\end{lstlisting}

\begin{retenirbox}[L'avantage des variables]
Si demain vous décidez de changer le vert pour un \texttt{\#1A9A5F} (un
peu plus clair), vous changez \textbf{une seule ligne} dans \texttt{:root}
et toute l'application est mise à jour. Sans variables, il faudrait
chercher et remplacer 20 occurrences à la main.
\end{retenirbox}

\begin{infobox}[\texttt{:root}, c'est quoi ?]
\texttt{:root} est un sélecteur spécial qui cible la racine du document
(c'est-à-dire \texttt{<html>}). On l'utilise pour déclarer des variables
\og{}globales\fg{} accessibles partout dans le CSS.
\end{infobox}

\newpage

% ===========================================================================
\section{Étape 4 : Importer les polices Google Fonts}
% ===========================================================================

Les polices par défaut du navigateur sont moches. \textbf{Google Fonts}
propose plus de 1000 polices gratuites, libres d'utilisation.

\subsection{Choisir 2 polices}

On va utiliser :
\begin{itemize}[nosep]
    \item \textbf{Poppins} pour le corps du texte (claire, ronde, moderne)
    \item \textbf{Montserrat} pour les titres (forte, bien lisible en gros)
\end{itemize}

\subsection{Méthode 1 : via le \texttt{<link>} HTML (recommandé)}

Dans le \texttt{<head>} de \textbf{chaque} page HTML, \textbf{avant} la
ligne \texttt{<link rel="stylesheet" href="style.css">}, ajoutez :

\begin{lstlisting}[language=HTML5]
<link rel="preconnect" href="https://fonts.googleapis.com">
<link rel="preconnect" href="https://fonts.gstatic.com" crossorigin>
<link href="https://fonts.googleapis.com/css2?family=Poppins:wght@400;600&family=Montserrat:wght@600;700&display=swap" rel="stylesheet">
\end{lstlisting}

\subsection{Utiliser les polices dans style.css}

\textbf{Mettez à jour le \texttt{body}} dans votre CSS :

\begin{lstlisting}[language=CSS]
body {
    font-family: 'Poppins', sans-serif;
    line-height: 1.6;
    color: var(--noir);
    background-color: var(--blanc-casse);
}

h1, h2, h3 {
    font-family: 'Montserrat', sans-serif;
    font-weight: 700;
}
\end{lstlisting}

\textbf{Enregistrez}. La page utilise maintenant Poppins pour le texte
courant et Montserrat pour les titres. \cmark

\begin{attentionbox}[Connexion internet obligatoire]
Google Fonts est servi en ligne. Sans Internet, les polices ne se chargent
pas et le navigateur retombe sur \texttt{sans-serif} par défaut. C'est
\textbf{normal} et ce n'est pas un bug. En production sérieuse, on
téléchargerait les fichiers de police pour les servir localement, mais
pour ce cours on garde la simplicité.
\end{attentionbox}

\newpage

% ===========================================================================
\section{Étape 5 : Styliser le header}
% ===========================================================================

\textbf{Ajoutez à \texttt{style.css}} :

\begin{lstlisting}[language=CSS]
/* ===== Header ===== */
header {
    background-color: var(--vert);
    color: white;
    padding: 30px 20px;
    text-align: center;
}

header h1 {
    font-size: 2.5rem;
    margin-bottom: 8px;
    letter-spacing: 1px;
}

header p {
    font-size: 1.1rem;
    opacity: 0.9;
    font-style: italic;
}
\end{lstlisting}

\textbf{Enregistrez}. Le bandeau du haut est maintenant vert vif avec un
gros titre blanc «\,Jàng Afrig\,» et un sous-titre italique en dessous. \cmark

\begin{infobox}[Quelques propriétés nouvelles]
\begin{itemize}[leftmargin=*,nosep]
    \item \texttt{padding: 30px 20px} \quad raccourci pour 30px haut/bas et 20px gauche/droite
    \item \texttt{font-size: 2.5rem} \quad 2.5 fois la taille de base ($\approx$ 40px)
    \item \texttt{letter-spacing: 1px} \quad espacement entre les lettres
    \item \texttt{opacity: 0.9} \quad transparence à 90\% (1 = opaque, 0 = invisible)
\end{itemize}
\end{infobox}

% ===========================================================================
\section{Étape 6 : Styliser la nav avec Flexbox}
% ===========================================================================

\textbf{Ajoutez à \texttt{style.css}} :

\begin{lstlisting}[language=CSS]
/* ===== Navigation ===== */
nav {
    background-color: var(--jaune);
    padding: 0;
    display: flex;
    justify-content: center;
    flex-wrap: wrap;
}

nav a {
    color: var(--noir);
    text-decoration: none;
    font-weight: 600;
    padding: 15px 25px;
    transition: background-color 0.2s;
}

nav a:hover {
    background-color: var(--vert);
    color: white;
}
\end{lstlisting}

\textbf{Enregistrez}. La nav est maintenant un bandeau jaune avec les
liens centrés en gras noir. \textbf{Survolez un lien} : il devient vert
avec texte blanc, en douceur. \cmark

\begin{retenirbox}[L'effet \texttt{:hover} avec \texttt{transition}]
La propriété \texttt{transition: background-color 0.2s} fait que le
changement de couleur au survol se fait \textbf{progressivement} sur
0.2 seconde au lieu d'être brutal. Petit détail, gros effet de
qualité perçue.
\end{retenirbox}

\subsection{Pourquoi le « | » a disparu ?}

Au Lab 2, la nav contenait des séparateurs \texttt{|} entre les liens.
Maintenant qu'on utilise Flexbox + padding, on n'en a plus besoin :
chaque lien a son propre espace. \textbf{Vous pouvez retirer les
\texttt{|}} dans le HTML des 5 pages. La nav devient :

\begin{lstlisting}[language=HTML5]
<nav>
    <a href="index.html">Accueil</a>
    <a href="quiz.html">Jouer</a>
    <a href="scores.html">Mes Scores</a>
    <a href="apropos.html">À propos</a>
</nav>
\end{lstlisting}

\newpage

% ===========================================================================
\section{Étape 7 : Cartes pour les modes de jeu}
% ===========================================================================

Les modes de jeu sont actuellement une simple liste \texttt{<ul>}. On va
les transformer en \textbf{cartes} colorées disposées en grille avec
Flexbox.

\subsection{Modifier le HTML de index.html}

\textbf{Remplacez} la \texttt{<section>} des modes de jeu dans
\texttt{index.html} par :

\begin{lstlisting}[language=HTML5]
<section>
    <h2>Les modes de jeu</h2>
    <div class="modes-grille">
        <div class="carte-mode carte-jaune">
            <h3>Capitales</h3>
            <p>« Quelle est la capitale du Mali ? »</p>
        </div>
        <div class="carte-mode carte-rouge">
            <h3>Drapeaux</h3>
            <p>« À quel pays appartient ce drapeau ? »</p>
        </div>
        <div class="carte-mode carte-vert">
            <h3>Langues</h3>
            <p>« Dans quel pays parle-t-on le swahili ? »</p>
        </div>
        <div class="carte-mode carte-jaune">
            <h3>Voisins</h3>
            <p>« Combien de voisins a l'Algérie ? »</p>
        </div>
        <div class="carte-mode carte-rouge">
            <h3>Chrono</h3>
            <p>Maximum de bonnes réponses en 60 secondes</p>
        </div>
    </div>
</section>
\end{lstlisting}

\subsection{Ajouter le CSS}

\textbf{Ajoutez à \texttt{style.css}} :

\begin{lstlisting}[language=CSS]
/* ===== Section principale ===== */
main {
    max-width: 1100px;
    margin: 30px auto;
    padding: 0 20px;
}

main section {
    margin-bottom: 40px;
}

main h2 {
    color: var(--vert);
    border-bottom: 3px solid var(--jaune);
    padding-bottom: 8px;
    margin-bottom: 20px;
}

/* ===== Grille de cartes (Flexbox) ===== */
.modes-grille {
    display: flex;
    flex-wrap: wrap;
    gap: 15px;
}

.carte-mode {
    flex: 1 1 200px;
    padding: 20px;
    border-radius: 8px;
    color: var(--noir);
    box-shadow: 0 2px 6px rgba(0, 0, 0, 0.1);
}

.carte-mode h3 {
    margin-bottom: 8px;
}

.carte-jaune  { background-color: var(--jaune); }
.carte-rouge  { background-color: var(--rouge); color: white; }
.carte-vert   { background-color: var(--vert);  color: white; }
\end{lstlisting}

\textbf{Enregistrez}. Les modes de jeu apparaissent maintenant en
\textbf{cartes colorées en grille}, qui se réorganisent automatiquement
si la fenêtre rétrécit. \cmark

\begin{infobox}[\texttt{flex: 1 1 200px}, c'est quoi ?]
Cette syntaxe condensée signifie : \og{}chaque carte peut grandir et
rétrécir, mais sa largeur \emph{idéale} est 200px\fg{}. Combiné avec
\texttt{flex-wrap: wrap} sur le parent, les cartes se mettent
automatiquement sur plusieurs lignes si l'écran est trop étroit. C'est
le début du \emph{responsive}.
\end{infobox}

\newpage

% ===========================================================================
\section{Étape 8 : Boutons stylés}
% ===========================================================================

Le bouton «\,Commencer\,» de la page \texttt{quiz.html} est moche. On
le transforme en bouton vert avec effet de survol.

\subsection{Ajouter une classe au bouton}

Dans \texttt{quiz.html}, modifiez le bouton :

\begin{lstlisting}[language=HTML5]
<button type="submit" class="btn btn-primary">Commencer le quiz</button>
\end{lstlisting}

\subsection{Styler les boutons en CSS}

\textbf{Ajoutez à \texttt{style.css}} :

\begin{lstlisting}[language=CSS]
/* ===== Boutons ===== */
.btn {
    display: inline-block;
    padding: 12px 28px;
    font-size: 1rem;
    font-weight: 600;
    border: none;
    border-radius: 6px;
    cursor: pointer;
    text-decoration: none;
    transition: background-color 0.2s, transform 0.1s;
}

.btn-primary {
    background-color: var(--vert);
    color: white;
}

.btn-primary:hover {
    background-color: #0a6e3f;  /* vert plus foncé */
}

.btn-primary:active {
    transform: scale(0.97);
}

.btn-secondary {
    background-color: var(--blanc-casse);
    color: var(--noir);
    border: 2px solid var(--noir);
}

.btn-secondary:hover {
    background-color: var(--noir);
    color: white;
}
\end{lstlisting}

\textbf{Enregistrez}. Le bouton est maintenant vert avec coins arrondis.
\textbf{Survolez-le} : il devient vert foncé. \textbf{Cliquez-le} : il
rétrécit légèrement (effet \og{}enfoncé\fg{}). \cmark

\begin{retenirbox}[Pseudo-classes pour les boutons]
\begin{itemize}[nosep]
    \item \texttt{:hover} \quad au survol de la souris
    \item \texttt{:active} \quad pendant le clic (souris enfoncée)
    \item \texttt{:focus} \quad quand sélectionné au clavier (\texttt{Tab})
\end{itemize}
Pour l'accessibilité, traitez \texttt{:hover} et \texttt{:focus} de la
même façon (pour les utilisateurs au clavier).
\end{retenirbox}

\newpage

% ===========================================================================
\section{Étape 9 : Tableau du palmarès}
% ===========================================================================

Le tableau de \texttt{scores.html} est sans aucune mise en forme. On lui
donne un look propre.

\textbf{Ajoutez à \texttt{style.css}} :

\begin{lstlisting}[language=CSS]
/* ===== Tableaux ===== */
table {
    width: 100%;
    border-collapse: collapse;
    margin-top: 15px;
    box-shadow: 0 2px 6px rgba(0, 0, 0, 0.05);
}

th {
    background-color: var(--vert);
    color: white;
    padding: 12px;
    text-align: left;
}

td {
    padding: 10px 12px;
    border-bottom: 1px solid #e0e0e0;
}

/* Lignes alternées (zébrures) */
tr:nth-child(even) td {
    background-color: rgba(245, 197, 24, 0.08);
}

tr:hover td {
    background-color: rgba(14, 138, 79, 0.1);
}
\end{lstlisting}

\textbf{Enregistrez} et allez sur \texttt{scores.html}. Le tableau est
maintenant : en-tête vert, lignes claires alternées avec un soupçon de
jaune, surlignage vert au survol. \cmark

\begin{infobox}[\texttt{tr:nth-child(even)}]
\texttt{:nth-child(even)} sélectionne les éléments \emph{pairs} (la 2\textsuperscript{e}, la 4\textsuperscript{e}, la 6\textsuperscript{e}\ldots). Combiné avec une couleur très claire, ça crée l'effet \og{}zébrures\fg{} qui aide à suivre les lignes du tableau.
\end{infobox}

% ===========================================================================
\section{Étape 10 : Footer}
% ===========================================================================

\textbf{Ajoutez à \texttt{style.css}} :

\begin{lstlisting}[language=CSS]
/* ===== Footer ===== */
footer {
    background-color: var(--noir);
    color: var(--blanc-casse);
    padding: 25px 20px;
    text-align: center;
    margin-top: 50px;
    font-size: 0.9rem;
}

footer a {
    color: var(--jaune);
    text-decoration: underline;
}

footer a:hover {
    color: white;
}
\end{lstlisting}

Le pied de page devient un bandeau noir avec texte clair et liens jaunes. \cmark

\newpage

% ===========================================================================
\section{Vérification visuelle}
% ===========================================================================

\subsection{Naviguez entre les 5 pages}

Avec Live Server, parcourez :
\begin{enumerate}[nosep]
    \item \texttt{index.html} : header vert, nav jaune, cartes colorées, footer noir
    \item \texttt{quiz.html} : formulaire avec bouton vert «\,Commencer\,»
    \item \texttt{resultat.html} : page de résultat
    \item \texttt{scores.html} : tableau du palmarès stylé
    \item \texttt{apropos.html} : sections informatives
\end{enumerate}

\textbf{Toutes les pages partagent le même look.} C'est la magie d'un
fichier CSS externe lié à toutes les pages. \cmark

\subsection{DevTools : modifier le CSS en direct}

\begin{infobox}[Inspecter et modifier]
Dans le navigateur :
\begin{enumerate}[nosep]
    \item Clic droit sur un élément $\to$ \textbf{Inspecter}
    \item L'onglet \textbf{Styles} (panneau de droite) montre toutes les règles CSS qui s'appliquent
    \item \textbf{Modifiez une valeur} (par exemple changez \texttt{background-color: var(--vert)} en \texttt{red}) : la page change instantanément
    \item Vos changements sont \textbf{temporaires} (ils disparaissent au rafraîchissement). C'est parfait pour expérimenter sans toucher au code.
\end{enumerate}
C'est l'outil que tous les développeurs utilisent pour bricoler.
\end{infobox}

\subsection{Ce qui ne marche pas encore (et c'est normal)}

\begin{itemize}[leftmargin=*,nosep]
    \item Le formulaire ne fait toujours rien quand on clique \og{}Commencer\fg{} → c'est pour le \textbf{Lab 7} avec JS
    \item L'app n'est pas encore optimisée pour téléphone → c'est pour le \textbf{Lab 4} (responsive)
    \item Les images ne sont pas dans une grille → c'est pour le \textbf{Lab 4} aussi
\end{itemize}

\newpage

% ===========================================================================
\section{Exercices d'application}
% ===========================================================================

\begin{objectifbox}[Trois mini-exercices à enchaîner]

\textbf{Exercice 1 --- Bouton secondaire}

Sur \texttt{resultat.html}, transformez les liens
\og{}Rejouer\fg{}, \og{}Voir mes scores\fg{}, \og{}Retour à l'accueil\fg{}
en \textbf{boutons} en utilisant les classes \texttt{.btn} et
\texttt{.btn-secondary} (déjà dans votre CSS) :

\begin{lstlisting}[language=HTML5]
<a href="quiz.html" class="btn btn-secondary">Rejouer</a>
<a href="scores.html" class="btn btn-secondary">Voir mes scores</a>
<a href="index.html" class="btn btn-secondary">Retour à l'accueil</a>
\end{lstlisting}

Bonus : disposez-les en ligne avec \texttt{display: flex; gap: 10px;} sur
leur conteneur.

\medskip

\textbf{Exercice 2 --- Personnaliser une carte de mode}

Modifiez la carte \og{}Chrono\fg{} sur \texttt{index.html} pour :
\begin{itemize}[nosep]
    \item Lui donner une nouvelle classe \texttt{.carte-noire}
    \item Définir cette classe en CSS : fond noir, texte blanc, bordure jaune épaisse
\end{itemize}

\medskip

\textbf{Exercice 3 --- Inverser les couleurs au survol des cartes}

Faites en sorte qu'au survol d'une carte de mode, ses couleurs
s'inversent (texte sur fond clair devient fond foncé sur texte clair).
Indice : utilisez \texttt{.carte-mode:hover} et \texttt{transition}.

\medskip

\textbf{Bonus} : changez la palette ! Remplacez le vert par une autre
couleur dominante (par exemple un rouge brique \texttt{\#a04436}) et
observez que toute l'app se met à jour grâce aux variables CSS.
\end{objectifbox}

\newpage

% ===========================================================================
\section{Aide-mémoire}
% ===========================================================================

\subsection{Propriétés CSS vues dans ce Lab}

\begin{center}
\begin{tabular}{@{}lp{8cm}@{}}
\toprule
\textbf{Propriété}              & \textbf{Rôle} \\
\midrule
\texttt{background-color}       & Couleur de fond \\
\texttt{color}                  & Couleur du texte \\
\texttt{font-family}            & Police de caractères \\
\texttt{font-size}              & Taille de police (px, rem, em) \\
\texttt{font-weight}            & Graisse (400 normal, 600 demi-gras, 700 gras) \\
\texttt{text-align}             & Alignement texte (\texttt{left}, \texttt{center}, \texttt{right}) \\
\texttt{text-decoration}        & Soulignement (\texttt{underline}, \texttt{none}) \\
\texttt{padding} / \texttt{margin} & Espacement intérieur / extérieur \\
\texttt{border} / \texttt{border-radius} & Bordure / arrondi des coins \\
\texttt{box-shadow}             & Ombre portée \\
\texttt{display: flex}          & Active Flexbox sur le conteneur \\
\texttt{flex} (raccourci)       & Comportement d'un flex item \\
\texttt{gap}                    & Espace entre les flex items \\
\texttt{transition}             & Animation douce sur changement de propriété \\
\texttt{cursor: pointer}        & Curseur en main au survol \\
\texttt{opacity}                & Transparence (0 à 1) \\
\bottomrule
\end{tabular}
\end{center}

\subsection{Variables et pseudo-classes}

\begin{itemize}[leftmargin=*,nosep]
    \item Définir une variable : \texttt{:root \{ --maCouleur: \#0E8A4F; \}}
    \item Utiliser une variable : \texttt{color: var(--maCouleur);}
    \item \texttt{:hover} \quad au survol
    \item \texttt{:active} \quad pendant le clic
    \item \texttt{:focus} \quad quand sélectionné
    \item \texttt{:nth-child(even)} \quad lignes paires
\end{itemize}

\subsection{Structure type d'un fichier style.css}

\begin{lstlisting}[language=CSS]
/* 1. Reset minimal     */
* { box-sizing: border-box; margin: 0; padding: 0; }

/* 2. Variables         */
:root { --vert: #0E8A4F; }

/* 3. Body de base      */
body { font-family: 'Poppins', sans-serif; }

/* 4. Sélecteurs d'éléments (header, nav, main, footer) */
header { ... }
nav { ... }

/* 5. Sélecteurs de classes (composants réutilisables) */
.btn { ... }
.carte-mode { ... }

/* 6. États (hover, focus, active) */
.btn:hover { ... }
\end{lstlisting}

\newpage

% ===========================================================================
\section{Routine Git de fin de Lab}
% ===========================================================================

Dans le terminal intégré de VS Code :

\begin{lstlisting}
git status
git add .
git commit -m "v0.3 : CSS panafricain + Google Fonts"
git push
git tag v0.3
git push --tags
\end{lstlisting}

Vérifiez sur \url{https://github.com/votre-username/jang-afrig} :
\begin{itemize}[nosep]
    \item Le fichier \texttt{style.css} est bien là
    \item Les 5 pages HTML contiennent le \texttt{<link>}
    \item Le tag \texttt{v0.3} apparaît dans l'onglet \textbf{Tags}
\end{itemize}

% ===========================================================================
\section{Rapport à produire}
% ===========================================================================

\begin{attentionbox}[Rapport du Lab 3 --- obligatoire]
\textbf{Format :} 1-2 pages, déposé dans \texttt{rapports/rapport-lab-3-NOM-Prenom.pdf}.

\medskip

\textbf{Contenu (4 sections) :}
\begin{enumerate}[nosep]
    \item \textbf{Ce que j'ai fait} : style.css créé et lié, palette définie, polices importées, header/nav/cartes/boutons/tableau/footer stylés, exercices d'application complétés.
    \item \textbf{Difficultés rencontrées} : variables qui ne fonctionnent pas, polices qui ne se chargent pas, problèmes de spécificité, conflits de styles\ldots
    \item \textbf{Comment j'ai résolu} : DevTools, MDN, aide d'un camarade.
    \item \textbf{Lien vers mon repo avec le tag v0.3} : \\
          ex. \url{https://github.com/fatou-diop/jang-afrig/releases/tag/v0.3}
\end{enumerate}

\medskip

\textbf{Bonus dans le rapport (facultatif) :} une capture d'écran de votre
page d'accueil stylée. C'est le premier rendu visuel de votre projet, ça
mérite d'être archivé.
\end{attentionbox}

\newpage

% ===========================================================================
\section{Résumé et transition}
% ===========================================================================

\subsection{Ce que vous avez accompli}

\begin{itemize}[leftmargin=*,nosep]
    \item Créé un fichier \texttt{style.css} unique pour les 5 pages \cmark
    \item Mis en place reset CSS et \texttt{box-sizing: border-box} \cmark
    \item Importé 2 polices Google Fonts \cmark
    \item Défini une palette panafricaine en variables CSS \cmark
    \item Stylisé header, nav, cartes, boutons, tableau, footer \cmark
    \item Ajouté des effets \texttt{:hover} et \texttt{transition} \cmark
    \item Posé le tag \texttt{v0.3} \cmark
\end{itemize}

\subsection{Ce qu'on fera au Lab 4}

\begin{itemize}[leftmargin=*,nosep]
    \item Faire en sorte que l'app soit \textbf{lisible sur téléphone}
    \item Découvrir les \textbf{media queries}
    \item Construire un \textbf{layout responsive} (1 colonne mobile, 2-3 colonnes desktop)
    \item Gérer une grille d'images (drapeaux côte à côte au lieu de empilés)
    \item Tag \texttt{v0.4}
\end{itemize}

\begin{projetbox}[À ce stade du semestre]
Votre application a maintenant un \textbf{vrai look professionnel sur
ordinateur}. La structure HTML est solide (Lab 2), le style est appliqué
(Lab 3). Manque le \textbf{responsive} (Lab 4) avant qu'on attaque la
partie interactivité (JavaScript).
\end{projetbox}

\vspace{0.5cm}

\begin{center}
\fbox{
\begin{minipage}{0.85\textwidth}
\centering
\vspace{0.5cm}
\textbf{\Large Lab 3 --- Terminé \cmark}\\[0.3cm]
\normalsize
Vous avez transformé un document brut en application stylée.\\
Place au responsive : que ça marche aussi sur téléphone.\\[0.2cm]
\textit{Rendez-vous au Lab 4 !}\\[0.3cm]
\vspace{0.3cm}
\end{minipage}
}
\end{center}

\end{document}
